\chapter{Considerações finais}
    \par Este trabalho de conclusão de curso propôs-se a investigar e demonstrar, por meio de um estudo de caso prático, como os princípios da Arquitetura Limpa podem ser implementados em uma API REST desenvolvida com o framework Laravel, partindo de sua estrutura nativa baseada no padrão Model-View-Controller (MVC). Notou-se uma carência de guias detalhados que traduzissem os conceitos teóricos de Robert C. Martin para o contexto específico de um dos mais populares frameworks PHP da atualidade.

    \par A hipótese central, de que o padrão MVC não precisaria ser descartado, mas sim reposicionado dentro da Arquitetura Limpa, foi confirmada pela refatoração da API de contatos. O estudo demonstrou que os componentes do MVC, como o Controller, foram adaptados para operar como detalhes na camada de Adaptadores de Interface, servindo como porta de entrada para o núcleo da aplicação, mas sem conter a lógica de negócio. Essa abordagem permitiu que a aplicação se beneficiasse da organização da Arquitetura Limpa sem abandonar completamente os padrões estabelecidos pelo framework.

    \par A execução do projeto evidenciou duas transformações fundamentais. A primeira foi a mudança estrutural, que evoluiu de uma organização por tipo técnico (Models, Controllers) para uma "Arquitetura Gritante", organizada por funcionalidades (UseCases/CreateContact), que comunica o propósito do sistema antes de revelar suas tecnologias. A segunda foi a alteração no fluxo de requisição, que passou de um modelo direto e fortemente acoplado para um fluxo desacoplado, orquestrado por interfaces e DTOs, que garante a estrita obediência à Regra da Dependência.

    \par Os resultados apontam que a adoção da Arquitetura Limpa introduz um trade-off claro. Por um lado, o principal ganho reside na manutenibilidade e flexibilidade a longo prazo, sendo este o objetivo central da arquitetura. O isolamento completo do núcleo de negócio - composto pelas Entidades e Casos de Uso - dos detalhes de infraestrutura, como o Eloquent ORM ou o próprio Laravel, protege a lógica mais valiosa do software contra mudanças tecnológicas e facilita sua evolução. 
    
    \par Por outro lado, um problema notado foi o aumento expressivo na quantidade de arquivos e na complexidade inicial. A criação de uma única funcionalidade, que antes se resumia a um método, passou a exigir a definição de múltiplas classes, como Entidades, DTOs, Boundaries e Interactors. Essa estrutura mais elaborada pode representar um custo de desenvolvimento inicial mais alto, e talvez não seja a solução mais eficiente para projetos pequenos ou com escopo muito simples, onde a agilidade imediata é mais crítica que a longevidade arquitetural.

    \par Em suma, este trabalho cumpre seu objetivo ao oferecer um guia prático e detalhado da refatoração, servindo como uma referência para desenvolvedores que buscam construir aplicações Laravel mais robustas. Conclui-se que a Arquitetura Limpa, mais do que um padrão prescritivo, é uma filosofia de design que representa um investimento estratégico na longevidade e na qualidade de um projeto de software.